% !TeX root = ../Book1.tex
\chapter*{序言}
\addcontentsline{toc}{chapter}{序言}
同一个圆，可以由距离方程描述，也可以看成一条闭合曲线、一个拓扑空间或一个群。每一种描述都保留某些性质，并提出不同的问题：哪些变换保持距离？沿圆周走一圈留下什么信息？局部参数怎样拼成整体？《几何学》从这些具体问题出发，逐步建立研究几何结构、空间与映射的方法。本卷讨论古典几何与光滑流形，是全书五卷中的第一卷。

\ProseParagraphBreak

古典几何中的坐标、变换和分类，为后来的抽象构造提供了可以反复检验的例子。圆锥曲线的标准形取决于允许仿射变换、等距变换还是射影变换；曲面的局部度量可以写在一张坐标图上，曲面整体的拓扑却要通过粘合来认识。进入流形以后，这些问题仍在：换坐标时什么量保持意义，局部给出的数据何时能相容，怎样由相容的数据构造整体对象。本卷的安排围绕这些问题展开。

\ProseParagraphBreak

本书面向正在首次或第二次系统学习几何的数学系本科生和研究生。开始阅读需要有限维线性代数和微积分的基本训练：能够使用基、线性映射、对偶、行列式与内积，理解导数、积分、极限和连续性。曲线曲面与流形部分还要用到多元微积分；有关微分方程、拓扑和群作用的准备，在实际使用处交代。已经学过某个定理却忘记了条件的读者，可以先看该处的回顾和引用，再决定是否回读其证明。

\ProseParagraphBreak

前两编先讨论几何对象怎样分类、空间怎样构造。第一编从仿射与 Euclid 结构进入射影几何和凸性，把坐标化简与不变量联系起来；第二编建立拓扑、基本群和覆盖，继而通过切割与粘合研究紧曲面。这里的分类同时包含两项工作：证明给出的模型能够容纳所有对象，并找出区分不同模型的依据。这样的区分也将用于后面的曲线基本定理、平坦结构和齐性空间。

\ProseParagraphBreak

第三编研究曲线与曲面的微分几何。曲率、基本形式和联络逐步把局部计算联系起来，曲面的整体公式则把曲率积分接到拓扑。球面与双曲模型让读者比较不同度量下的距离、测地线和群作用。第四编将这些构造放到一般光滑流形上，依次处理切空间、向量丛、流、微分形式与联络。第五编研究连续对称性，借助 Lie 群、齐性空间和主丛，把换标架、群作用与平行移动纳入共同的几何描述；最后两章继续讨论紧群的结构与表示。

\ProseParagraphBreak

读者可以按自己的目的选择起点。希望从熟悉的图形逐步学习者，可以依次阅读前三编，再进入流形。已经具备线性代数和多元微积分基础、希望直接学习光滑流形者，可以从\chapref{g1:ch:11}开始；所需的开集、连续映射、同胚、子空间与积空间，可在\chapref{g1:ch:04}查阅。逆函数、隐函数与常秩定理在\chapref{g1:ch:11}重新建立，单位分解的光滑构造在\chapref{g1:ch:12}展开。读到平坦联络与单值时，再补读\chapref{g1:ch:05}的基本群和覆盖理论。古典分类与曲面分类可以另行阅读。

\ProseParagraphBreak

各编导读说明本编的问题和所需基础，各章的摘要与学习目标便于确认当前任务。正文中的例子既用于理解定义，也用于检验条件；同一个环面、射影空间或球面再次出现时，可以比较这一次新增了哪些结构。习题配有解答，阅读时宜先独立计算或构造，再用解答检查关键步骤。对于较长证明，先辨认待构造的对象及主要障碍，再跟随局部计算回到完整结论，通常比逐式记忆更有帮助。

\ProseParagraphBreak

章、节及习题中的星号标明可按学习目的选读的内容，附录提供补充专题。附录 A 讨论公理化几何与几何史的方法；附录 B、E 延伸离散、凸性与反射群问题；附录 C 讨论计算与图示，附录 D 继续研究几何结构。它们可以在相应正文之后选读。各章的阅读说明给出适合对照学习的教材，具体引文则标明所用结果或证明的出处；专题专著和原始论文用于继续追问某一问题。卷末参考文献与索引供检索，阅读时以当前问题所需的资料为先。

\ProseParagraphBreak

后续各卷分别讨论代数拓扑与微分拓扑、黎曼几何与几何分析、复几何与辛几何、代数几何。本卷的流形、丛、微分形式和联络会在其中反复使用，所需的拓扑、分析或代数工具则随问题继续建立。当前二十章、附录及书前书后文字均为初稿，尚未经过正式审核。
